\documentclass[12pt,a4paper]{article}
\usepackage[utf8]{inputenc}
\usepackage[spanish,es-noquoting]{babel}
\usepackage[T1]{fontenc}
\usepackage{lmodern}
\usepackage{geometry}
\geometry{top=2cm, bottom=2cm, left=2.5cm, right=2.5cm}
\usepackage{xcolor}
\usepackage{listings}
\usepackage{enumitem}
\usepackage{fancyhdr}
\usepackage{hyperref}
\usepackage{tikz}
\usetikzlibrary{arrows.meta, positioning, calc}

\definecolor{codegreen}{rgb}{0,0.6,0}
\definecolor{backcolour}{rgb}{0.95,0.95,0.92}

\lstset{
    language=C,
    backgroundcolor=\color{backcolour},
    commentstyle=\color{codegreen},
    keywordstyle=\color{blue},
    basicstyle=\ttfamily\small,
    breaklines=true,
    frame=single,
    numbers=left,
    tabsize=4,
    extendedchars=true,
    showstringspaces=false,
    inputencoding=utf8,
    literate={á}{{\'a}}1 {é}{{\'e}}1 {í}{{\'i}}1 {ó}{{\'o}}1 {ú}{{\'u}}1 {ñ}{{\~n}}1 {¡}{{\textexclamdown}}1 {¿}{{\textquestiondown}}1
}

\pagestyle{fancy}
\fancyhf{}
\lhead{Monitorías Sistemas Operativos 2026-I}
\rhead{Capítulo 3: Espera y Finalización}
\cfoot{\thepage}

\title{\textbf{Guía de Aprendizaje: Sincronización Básica (\texttt{pthread\_join})}}
\author{Malak Sanchez \\ \small Monitorías de Sistemas Operativos}
\date{\today}

\begin{document}

\maketitle

\section{Motivación}
Cuando creamos un hilo, este corre de forma paralela al hilo principal. Si el programa principal termina (llega al final de \texttt{main} o llama a \texttt{exit()}), el proceso completo muere, matando a todos los hilos en su interior, sin importar si han terminado su trabajo o no. Necesitamos una forma de esperar a que finalicen.

\section{Idea Fundamental: Join}
\texttt{pthread\_join} es el equivalente hilos a la función \texttt{wait()} de los procesos. Su función es bloquear al hilo que hace la llamada hasta que el hilo especificado termine su ejecución.

\begin{center}
\begin{tikzpicture}[font=\small, >=Stealth, 
    thread_path/.style={ultra thick, violet!70!black},
    main_path/.style={ultra thick, blue!60!black}]

    % Vertical timeline
    \draw[->, thick, gray, opacity=0.2] (0, 5) -- (0, 0) node[below] {Tiempo};

    % Main thread
    \draw[main_path] (-1, 5) -- (-1, 4) node[midway, left] {Main Thread};
    
    % Thread creation
    \draw[->, thick, dashed] (-1, 4) -- (1, 3.5) node[midway, above] {create()};
    
    % New thread execution
    \draw[thread_path] (1, 3.5) -- (1, 1.5) node[midway, right] {New Thread};
    
    % Main waiting
    \draw[main_path, gray!50] (-1, 4) -- (-1, 1.5) node[midway, left] {Blocked (Join)};
    
    % Completion and Join
    \draw[->, thick, dashed] (1, 1.5) -- (-1, 1.5) node[midway, below] {Fin / Join};
    
    % Main resumes
    \draw[main_path] (-1, 1.5) -- (-1, 0.5) node[midway, left] {Resume};

\end{tikzpicture}
\end{center}

\section{Sintaxis}
\begin{lstlisting}[numbers=none]
int pthread_join(pthread_t thread, void **retval);
\end{lstlisting}

\begin{itemize}
    \item \texttt{thread}: identificador del hilo que se desea esperar. 
    La función bloquea el hilo que la invoca hasta que este hilo termine su ejecución.

    \item \texttt{retval}: puntero donde se almacenará el valor de retorno del hilo
    cuando finalice. Si no se desea recuperar dicho valor, se puede pasar
    \texttt{NULL}.
\end{itemize}

\newpage

\section{Ejemplo: Retornando un Valor}
\begin{lstlisting}
#include <stdio.h>
#include <stdlib.h>
#include <pthread.h>

void* compute_result(void* arg) {
    int *result = (int*) malloc(sizeof(int));
    *result = 42;
    printf("Thread: Computation finished.\n");
    pthread_exit((void*)result); // Return the pointer
}

int main(){
    pthread_t thread;
    void* result_ptr;

    pthread_create(&thread, NULL, compute_result, NULL);

    // Stop and wait here for the thread
    pthread_join(thread, &result_ptr);

    printf("Main: Thread returned %d\n", *(int*)result_ptr);
    
    free(result_ptr);
    return 0;
}
\end{lstlisting}

\section{Hilos Detach (Desconectados)}
Si no necesitamos esperar a un hilo ni obtener su valor de retorno, podemos marcarlo como \textbf{detached}. Esto libera sus recursos automáticamente al terminar sin necesidad de un join.
\begin{lstlisting}[language=C, numbers=none]
pthread_detach(thread_id);
\end{lstlisting}

\section{Ejercicios}
\begin{enumerate}
    \item Cree tres hilos que realicen cálculos diferentes (sumas simples) y devuelvan el resultado al hilo principal usando \texttt{pthread\_join}.
    \item Investigue qué sucede si intenta hacer \texttt{pthread\_join} sobre un hilo que ya fue "unido" o que está en estado "detached".
\end{enumerate}

\end{document}
